\subsection{Алгоритм GraphRAG}

В качестве референсного решения был выбран ObsidianRAG~\cite{obsidianrag}.
Его алгоритм состоит из нескольких шагов:
\begin{enumerate}
      \item Индексация: содержимое \texttt{markdown}-файлов разбивается на чанки
            разметом 1500 символом с перекрытием в 300 символов.
            Полученные фрагменты вместе с извлеченными из текста ссылками
            записываются в ChromaDB - векторную базу данных -
            которая создает для них плотные векторные эмбеддинги.
      \item Первоначальный поиск: запрос пользователя обрабатывается
            с помощью \texttt{EnsembleRetriever}, использующим BM25 и ANN
            с весами 0.4 и 0.6 соответственно.
            Извлекается N наиболее релевантных заметок.
      \item Переранжирование: полученные файлы передаются в Cross Encoder,
            например, \texttt{BAAI/bge-reranker-v2-m3}.
            В результате фильтрации остается 6 наиболее близких фрагментов.
      \item Расширение контекста: результат дополняется пятью файлами,
            связанными ссылками с оставшимися.
      \item Генерация ответа: собранное содержимое заметок
            передается в локальную LLM (Ollama).
\end{enumerate}

У данного пайплайна есть ряд недостатков:

\begin{enumerate}
      \item Одноэтапное расширение контекста соседями.
            При таком подходе теряются сложные связи, требующие многошаговых рассуждений.
            А также возрастает вероятность остаться в семантической ловушке: ситуации,
            когда новая информация не может быть извлечена
            из-за сильной разницы в эмбеддингах и ключевых словах.
            Суть ссылок заключается в решении этой проблемы, но автор ими пренебрегает.
      \item Сбор контекста без учета его релевантности.
            В проекте уже используется Cross Encoder для первоначального поиска,
            его можно легко переиспользовать для фильтрации фрагментов, не относящихся к запросу.
\end{enumerate}

На основе этого решения был разработан улучшенный четырехэтапный пайплайн,
поддерживающий типизированные ссылки и многошаговое расширение контекста.
Он нацелен на решение фундаментальных проблем эталонного решения.
Алгоритм реализован на языке Python
(исходный код доступен в папке \texttt{src/graphrag}).

\begin{figure}[H]
      \centering
      \begin{tikzpicture}[
                  node distance = 6mm and 32mm,
                  block/.style = {
                              rectangle, rounded corners=3pt,
                              draw=black!60, fill=blue!8,
                              text width=76mm, minimum height=9mm,
                              align=center, font=\small
                        },
                  ioblock/.style = {
                              rectangle, rounded corners=3pt,
                              draw=black!40, fill=gray!10,
                              text width=76mm, minimum height=8mm,
                              align=center, font=\small\itshape
                        },
                  sideblock/.style = {
                              rectangle, rounded corners=3pt,
                              draw=black!50, fill=orange!12,
                              text width=38mm, minimum height=8mm,
                              align=center, font=\small
                        },
                  arrow/.style  = {-Stealth, thick, black!70},
                  dasharrow/.style = {dashed, thick, black!50, -Stealth},
                  label/.style  = {font=\scriptsize\itshape, text=black!60}
            ]

            % Main flow
            \node[ioblock]  (query)   {Запрос пользователя};
            \node[block, below=of query]   (seed)
            {1. Первоначальный поиск\\
            {\scriptsize\normalfont Гибридный поиск (BM25~+~вектор)}};
            \node[block, below=of seed]    (trav)
            {2. Обход соседей\\
            {\scriptsize\normalfont базовый или типизированный BFS}};
            \node[block, below=of trav]    (prune)
            {3. Отбор контекста\\
            {\scriptsize\normalfont Чанкинг, переранжирование}};
            \node[block, below=of prune]   (gen)
            {4. Генерация ответа\\
            {\scriptsize\normalfont LLM, структурированный промпт}};
            \node[ioblock, below=of gen]   (answer)  {Ответ};

            % Side block
            \node[sideblock, right=of trav] (kg)
            {Граф знаний\\{\scriptsize \texttt{rel\_type:: [[T]]}}};

            % Main arrows
            \draw[arrow] (query)  -- node[label, right]{} (seed);
            \draw[arrow] (seed)   -- node[label, right]{TOP\_K\_SEED заметок} (trav);
            \draw[arrow] (trav)   -- node[label, right]{заметки, дополненные соседями} (prune);
            \draw[arrow] (prune)  -- node[label, right]{TOP\_K\_CONTEXT релевантных чанков} (gen);
            \draw[arrow] (gen)    -- (answer);

            % Dashed arrow from KG to traversal
            \draw[dasharrow] (kg.west) -- node[label, above]{\scriptsize исходящие ссылки} (trav.east);

      \end{tikzpicture}
      \caption{Общая схема конвейера GraphRAG}
      \label{fig:graphrag_pipeline}
\end{figure}

\textbf{Этап 1. Первоначальный поиск.}
На этом этапе выполняется гибридный поиск по LanceDB.
Он включает в себя ANN-поиск по плотным эмбеддингам и BM25-поиск по
лемматизированному тексту. На его основе собирается $TOP\_K\_SEED$ (по умолчанию~$= 5$)
наиболее близких заметок-кандидатов. Они добавляются в множество посещенных файлов.

В отличие от референсного решения,
на данном этапе не выполняется переранжирование заметок с помощью Cross Encoder.
Оно было перенесено на шаг 3.
Это решение было принято исходя из точности алгоритма извлечения, а также для
смягчения критериев для отбора заметок,
что потенциально может увеличить полноту собранного контекста.

\textbf{Этап 2. Обход соседей.}

Расширение контекста может быть проведено двумя способами:

\begin{itemize}
      \item Базовый обход в ширину. Для каждой заметки формируется список исходящих ссылок.
            Типы ссылок при этом игнорируются, как и релевантность их содержимого.
            Все непосещенные файлы, соединенные с текущим, добавляются в очередь для обхода.
      \item Типизированный обход в ширину. Для каждой заметки, связанной с текущей,
            формируется запрос в Cross Encoder модель формата:
            \begin{center}
                  \texttt{relation: target. summary}
            \end{center}
            где \texttt{relation}~--- тип связи, \texttt{target}~--- имя целевой заметки,
            \texttt{summary}~--- ее краткое содержание (первый абзац).
            Cross Encoder вычисляет оценку близости соседней заметки с текущей.
            Все кандидаты с оценкой ниже пороговой не добавляются в очередь обхода.
            Это позволяет фильтровать заметки на этапе сбора контекста, что
            может снизить количество шума, но также уменьшить полноту контекста из-за раннего отброса.
\end{itemize}

Результатом этапа является множество потенциально релевантных заметок.

Наглядное сравнение стратегий обхода представлено в таблице~\ref{tab:traversal_compare}.

\begin{table}[H]
      \centering
      \begin{tabularx}{\textwidth}{|l|X|X|}
            \hline
            \textbf{Критерий}             & \textbf{Base}    & \textbf{Typed}                     \\
            \hline
            Тип связи                     & Игнорируется     & Включается в оценку                \\
            \hline
            Оценка кандидатов             & Отсутствует      & Cross Encoder                      \\
            \hline
            Текст для оценки              & ---              & \texttt{relation: target. summary} \\
            \hline
            Фильтрация                    & Нет (все соседи) & Пороговая по оценке                \\
            \hline
            Ограничение числа новых узлов & Нет              & Нет                                \\
            \hline
      \end{tabularx}
      \caption{Сравнение стратегий обхода графа}
      \label{tab:traversal_compare}
\end{table}

\textbf{Этап 3. Отбор содержимого.}
Содержимое каждой заметки, отобранной на предыдущем этапе, разбивается на фрагменты
с помощью того же сплиттера, что используется при индексации хранилища (глава 2.1).

Каждый полученный чанк получается оценку Cross Encoder схожести с запросом пользователя.
Сохраняется $TOP\_K\_CONTEXT$ (по умолчанию 10) наиболее релевантных фрагментов,
которые сформируют итоговый контекст и будут отправлены в LLM.

В отличие от референсного решения, в котором фильтрация происходит на этапе сбора контекста,
здесь происходит отбор всего собранного содержимого заметок после обхода графа.
Это позволяет увеличить полноту выборки,
но допускает наличие в ней большего количества нерелевантных фрагментов.

\textbf{Этап 4. Генерация ответа.}
Запрос в LLM формируется на основе промпта, который дает модели необходимые инструкции
(например, отвечать на языке пользователя, несмотря на язык предоставленного контекста,
и опираться только на предоставленные данные),
сообщения пользователя и
отобранных фрагментов текстов \texttt{markdown}-файлов,
разграниченные специальным разделителем.

Предсказание модели преобразуется в  формат JSON с помощью
\texttt{JsonOutputParser} из LangChain.
В ответе ожидается два поля:
\texttt{reasoning}~--- рассуждения модели, и
\texttt{answer}~--- финальный ответ, возвращаемый пользователю.
В случае неудачного парсинга возвращается сырой ответ модели.
